\begin{filecontents*}[overwrite]{refs.bib}
% Entries marked "VERIFY" were filled from secondary notes. Check them against the primary source.

@misc{takehi2026wsp,
  title  = {Sparks of Weight-Space Programming},
  author = {Takehi, Rikiya and Teoh and Zhang and Khattab, Omar},
  year   = {2026},
  howpublished = {\url{https://rikiyat.github.io/blog/weight-space-programming/}},
  note   = {MIT CSAIL, 2 November 2026. VERIFY: full author names}
}

@inproceedings{steer2edit2026,
  title  = {Steer2Edit},
  author = {Sun and others},
  booktitle = {Advances in Neural Information Processing Systems},
  year   = {2026},
  note   = {arXiv:2602.09870. VERIFY: full title and authors}
}

@inproceedings{zhou2025safetyheads,
  title  = {On the Role of Attention Heads in Large Language Model Safety},
  author = {Zhou, Zhenhong and others},
  booktitle = {International Conference on Learning Representations},
  year   = {2025},
  note   = {arXiv:2410.13708. VERIFY: author list}
}

@inproceedings{wu2025retrievalheads,
  title  = {Retrieval Head Mechanistically Explains Long-Context Factuality},
  author = {Wu, Wenhao and Wang, Yizhong and Xiao, Guangxuan and Peng, Hao and Fu, Yao},
  booktitle = {International Conference on Learning Representations},
  year   = {2025},
  note   = {arXiv:2404.15574}
}

@inproceedings{arditi2024refusal,
  title  = {Refusal in Language Models Is Mediated by a Single Direction},
  author = {Arditi, Andy and Obeso, Oscar and Syed, Aaquib and Paleka, Daniel and Panickssery, Nina and Gurnee, Wes and Nanda, Neel},
  booktitle = {Advances in Neural Information Processing Systems},
  year   = {2024},
  note   = {arXiv:2406.11717}
}

@misc{lavi2025unanswerability,
  title  = {Detecting (Un)answerability in Large Language Models with Linear Directions},
  author = {Lavi and Milo and Geva, Mor},
  year   = {2025},
  note   = {arXiv:2509.22449. VERIFY: title and first names}
}

@inproceedings{lee2025cast,
  title  = {Programming Refusal with Conditional Activation Steering},
  author = {Lee, Bruce W. and others},
  booktitle = {International Conference on Learning Representations},
  year   = {2025},
  note   = {arXiv:2409.05907. VERIFY: author list}
}

@inproceedings{joren2025sufficient,
  title  = {Sufficient Context: A New Lens on Retrieval Augmented Generation Systems},
  author = {Joren, Hailey and others},
  booktitle = {International Conference on Learning Representations},
  year   = {2025},
  note   = {arXiv:2411.06037. VERIFY: author list}
}

@misc{kirichenko2025abstentionbench,
  title  = {AbstentionBench: Reasoning LLMs Fail on Unanswerable Questions},
  author = {Kirichenko, Polina and others},
  year   = {2025},
  note   = {arXiv:2506.09038. VERIFY}
}

@article{kumaran2026confidence,
  title  = {Confidence and abstention in large language models},
  author = {Kumaran, Dharshan and Daw and Osindero, Simon and Veli{\v{c}}kovi{\'c}, Petar and P{\u{a}}tr{\u{a}}ucean, Viorica},
  journal = {Nature Machine Intelligence},
  year   = {2026},
  note   = {arXiv:2603.22161. VERIFY: title and first names}
}

@inproceedings{shaham2024maia,
  title  = {A Multimodal Automated Interpretability Agent},
  author = {Shaham, Tamar Rott and others},
  booktitle = {International Conference on Machine Learning},
  year   = {2024},
  note   = {arXiv:2404.14394}
}

@misc{barez2026agenda,
  title  = {Automated interpretability research agenda},
  author = {Barez, Fazl},
  year   = {2026},
  note   = {University of Oxford. VERIFY: full citation}
}

@inproceedings{meng2022rome,
  title  = {Locating and Editing Factual Associations in {GPT}},
  author = {Meng, Kevin and Bau, David and Andonian, Alex and Belinkov, Yonatan},
  booktitle = {Advances in Neural Information Processing Systems},
  year   = {2022},
  note   = {arXiv:2202.05262}
}

@inproceedings{meng2023memit,
  title  = {Mass-Editing Memory in a Transformer},
  author = {Meng, Kevin and Sharma, Arnab Sen and Andonian, Alex and Belinkov, Yonatan and Bau, David},
  booktitle = {International Conference on Learning Representations},
  year   = {2023},
  note   = {arXiv:2210.07229}
}

@inproceedings{liu2026broken,
  title  = {Broken Symmetry in {LLM} Refusal: Answer Release Is More Local Than Refusal Restoration},
  author = {Liu and Wang and Wang and Xiao and Lin},
  booktitle = {Proceedings of EMNLP},
  year   = {2026},
  note   = {arXiv:2608.15772. VERIFY: first names}
}

@inproceedings{ling2025artifact,
  title  = {{LLM} Abstention Can Be a Prompt Artifact, in Addition to Genuine Uncertainty},
  author = {Ling and others},
  booktitle = {Proceedings of EMNLP},
  year   = {2026},
  note   = {arXiv:2507.16199. VERIFY: authors}
}

@misc{commitabstain2026,
  title  = {The Commit--Abstain Circuit},
  author = {{VERIFY}},
  year   = {2026},
  note   = {arXiv:2609.32964v2. VERIFY: title and authors}
}

@inproceedings{wang2025falserefusal,
  title  = {Surgical, Cheap, and Flexible: Mitigating False Refusal in Language Models via Single Vector Ablation},
  author = {Wang, Xinpeng and Hu, Chengzhi and R{\"o}ttger, Paul and Plank, Barbara},
  booktitle = {International Conference on Learning Representations},
  year   = {2025},
  note   = {arXiv:2410.03415. VERIFY: venue and author list}
}

@misc{hinkley2026paving,
  title  = {Residual Paving: Selective Refusal Editing},
  author = {Hinkley and Najafirad},
  year   = {2026},
  note   = {arXiv:2605.20262. VERIFY: first names and exact title}
}

@misc{liu2026auditors,
  title  = {Computation, not Weight Access: {LLMs} as Auditors of Clinical Prediction Networks},
  author = {Liu and Yang},
  year   = {2026},
  note   = {arXiv:2611.07847. VERIFY: first names}
}

@misc{hasanov2026drno,
  title  = {Dr. No},
  author = {Hasanov and others},
  year   = {2026},
  note   = {arXiv:2611.03862. VERIFY: full title and authors}
}

@inproceedings{setiawan2026grabrag,
  title  = {{GRAB-RAG}},
  author = {Setiawan},
  booktitle = {Proceedings of AACL},
  year   = {2026},
  note   = {arXiv:2608.22228. VERIFY: full title and first name}
}

@misc{readablenotsteerable2026,
  title  = {Readable Is Not Steerable},
  author = {{VERIFY}},
  year   = {2026},
  note   = {arXiv:2611.05577. Known from secondary coverage only. VERIFY: title and authors}
}
\end{filecontents*}

\documentclass[11pt]{article}

\usepackage[margin=1in]{geometry}
\usepackage[T1]{fontenc}
\usepackage{lmodern}
\usepackage{microtype}
\usepackage{amsmath,amssymb}
\usepackage{booktabs}
\usepackage{tabularx}
\usepackage{longtable}
\newcolumntype{Y}{>{\raggedright\arraybackslash}X}
\usepackage{graphicx}
\usepackage{xcolor}
\usepackage{newunicodechar}
\newunicodechar{§}{\S}
\usepackage[numbers,sort&compress]{natbib}
\usepackage[colorlinks=true,linkcolor=blue!50!black,citecolor=blue!50!black,urlcolor=blue!50!black]{hyperref}
\usepackage{cleveref}
\usepackage{enumitem}
\setlist{itemsep=2pt,topsep=4pt}

% Writing aids. Set \drafttrue to \draftfalse to hide every note before release.
\newif\ifdraft \drafttrue
\newcommand{\todo}[1]{\ifdraft\textcolor{red}{[TODO: #1]}\fi}
\newcommand{\pending}[1]{\ifdraft\textcolor{orange}{[PENDING #1]}\fi}
\newcommand{\tbd}{\textcolor{orange}{---}}   % empty result cell; never fill with a guess

\title{Evaluating Agent-Proposed Weight Edits for False Abstention:\\
\large A controlled case study in Qwen2.5-3B-Instruct}
\author{Qamil Mirza\\ \small\todo{affiliation, contact}}
\date{Technical report v0.2 (draft) --- \today}

\begin{document}
\maketitle

\begin{abstract}
Small instruction-tuned language models sometimes answer ``unknown'' even when their context contains enough
evidence to answer correctly. In Qwen2.5-3B-Instruct, this false abstention can follow a note in which the model
has already quoted the relevant fact. Earlier interventions recovered correct answers but also encouraged
unsupported answers. We ask whether an agent-proposed weight edit can improve this trade-off.
Following weight-space programming \citep{takehi2026wsp}, a Claude agent inspected the model on an exploration
split and proposed three edits. We evaluated them in a pre-registered case study in one model, with separate
selection and sealed confirmation splits, a transfer split, and six frozen comparators serving repair or pathway
tests. The selected repair was a rank-1 change to the attention output projection in layer 35. It met the
pre-registered repair criterion on sealed data: the base-weighted recovery rate on false abstentions was 82\%,
while unsupported answers on unanswerable questions rose from 3.5\% to 4.2\%. It damaged already-correct answers
at a base-weighted rate of 3\%. Its release-versus-unsupported-answer trade-off was better than a prompt
instruction, activation steering and directional orthogonalisation. The comparison with a simple logit bias
against ``unknown'' was inconclusive. A separate two-head edit provided evidence about a pathway through which
answers can be released, but the agent's explanation of that component was not supported by the pre-registered
tests.

\end{abstract}

\noindent\fbox{\parbox{\dimexpr\linewidth-2\fboxsep-2\fboxrule}{\small
\textbf{Key findings from the sealed evaluation}
\begin{itemize}\setlength\itemsep{1pt}
  \item \textbf{The selected edit met the repair criterion.} The rank-1 edit to the layer-35 attention output
        improved utility by $+0.117$ (95\% CI 0.096 to 0.138). It recovered the correct city on 212 of 261
        false abstentions. Base-weighted unsupported-answer rates rose from 0.035 to 0.042, and damage to
        already-correct answers was 0.030.
  \item \textbf{The comparison with simpler fixes was mixed.} The edit had a better trade-off than an added
        prompt instruction, a layer-29 steering intervention and directional orthogonalisation, after Holm
        correction. Superiority over removing the ``unknown'' option was not established. Against logit bias,
        neither superiority nor equivalence was established.
  \item \textbf{New-city results were consistent with confirmation.} On the descriptive transfer set, utility
        was 0.422 versus 0.303 for the unedited model, recovery was 0.817, and unsupported answers were 0.062
        versus 0.059. This was not a second confirmation test.
  \item \textbf{Repair and explanation were tested separately.} For a different, two-head edit, a layer-33
        intervention was sufficient but not necessary to release answers under the tested conditions.
        The agent's account of what the component computes was not supported by the frozen tests.
\end{itemize}}}
\vspace{0.5em}


\section{The question and main finding}
\label{sec:question}

A model receives a short memory stating where someone lives. Asked for the city, it writes a note quoting the
relevant fact, then answers ``unknown''. This is one form of \emph{false abstention}: the context supports an
answer, but the model declines to give it. An appropriate abstention is different---the context really does not
supply enough evidence. A useful intervention must recover the missed answers while preserving that distinction.
The note is a prompted summary, not evidence that the model completed a reasoning process
(\cref{sec:setup}).

This question grew out of a project on memory management. The broader story appears in
\href{https://qamil-mirza.com/research/more-than-reasoning}{\emph{More Than Reasoning}} and
\href{https://qamil-mirza.com/research/ai-brain-surgery}{\emph{Weight-Space Programming for Mechanistic Interpretability}}.
Here we focus on one experiment: can an AI agent propose a small weight edit that improves answering on examples
it has never seen, and does that edit offer a better trade-off than simpler interventions?

\paragraph{What motivated the experiment.}
Earlier experiments revealed a gap between recovering an answer-related signal and explaining a failure.
On Qwen2.5-3B-Instruct, the gold answer's first token appeared in the final-layer top-10 on 0.64 of failures;
the corresponding rates were 0.68--0.80 in three other model families. This is a readout-qualified observation,
not proof that the answer was derived from the supplied evidence. On controlled tasks, vocabulary lenses
recovered answers only in the last six layers. That timing reflects the readout's limits as well as the
model's computation.

The model's own generated notes made the problem especially visible. Of 165 held-out false abstentions,
152 named the gold city before answering ``unknown''. The city was written down; that alone does not show
that the model knew the answer. Asking the model whether the memory was sufficient did not resolve the issue:
Qwen 3B answered ``no'' on most rows it could answer correctly.

Instructions and activation steering could recover answers. A content-free direction added at layer 29
encouraged answering, following the refusal-direction line of work
\citep{arditi2024refusal,lavi2025unanswerability}. The evaluated ungated interventions also increased unsupported
answers on unanswerable questions: rates reached 0.188 for the layer-29 push and 0.451 for removing the
``unknown'' option. A separate, template-limited gating result qualified this pattern; it did not establish a
general solution. Some earlier observations replicated known results, and the steering work included
exploratory and secondary analyses.

False abstention was more frequent in Qwen 3B under the tested conditions: its in-view refusal rate was 0.30,
compared with 0.10 in Llama 3B, 0.02 in Gemma 4B and 0.03 in Phi-4-mini, which more often answered incorrectly
instead. This is a frequency contrast, not a claim that the behaviour is unique to Qwen.

\paragraph{Main finding.}
An agent-proposed edit to the final attention layer met the pre-registered repair criterion on sealed data.
Its trade-off between recovering correct answers and producing unsupported ones was better than an added prompt
instruction, an activation-steering intervention and directional orthogonalisation. Superiority over a logit bias
against ``unknown'', or over removing that option from the prompt, was not established. The following sections
describe what the agent could change, how the edits were selected, and what the independent tests support.

\section{What's new here}
\label{sec:new}

% The first-in-setting claim and the descriptions below still need a complete primary-source audit.
An agent proposing a useful weight edit is not itself new. Weight-space programming demonstrated that approach
\citep{takehi2026wsp}, and Steer2Edit studied edits to individual heads and neurons \citep{steer2edit2026}.
Earlier work also identified attention heads involved in safety refusal \citep{zhou2025safetyheads}.

Selectivity is an established concern. False-refusal vector ablation folded an intervention into model weights
and evaluated released false refusals against preserved true refusals \citep{wang2025falserefusal}.
Residual Paving used a gated residual update to address three objectives: change the target behaviour, preserve
benign behaviour, and retain wanted refusals \citep{hinkley2026paving}. Research on prompt-induced abstention
and local answer release also motivates testing late-layer interventions
\citep{liu2026broken,ling2025artifact}. Pre-registered, once-only held-out tests have precedent in
interpretability outside language models \citep{readablenotsteerable2026}.
\Cref{app:related} compares these studies with ours.

\paragraph{The contribution of this report.}
This is a pre-registered case study in one model. The proposed contribution is the combination of four elements:
false abstention with sufficient evidence in context; evaluation of answer recovery, unsupported answers, damage
to correct answers and following changed evidence; sealed confirmation of agent-proposed edits; and comparison
with frozen prompt, steering, orthogonalisation and logit-bias interventions. To our knowledge, this is the first
pre-registered test on sealed data of an agent-proposed weight edit for false abstention with the evidence in view.

The result has a clear boundary. The selected edit improved the pre-registered trade-off over three of the
simpler interventions. Against a logit bias on ``unknown'', however, the comparison was inconclusive. The point
estimates favoured the edit on recovery (0.816 versus 0.793) and unsupported answers (0.042 versus 0.095), but
the edit damaged more already-correct answers (0.030 versus none). The comparison did not survive correction
for multiple tests (Holm $p = 0.092$). We found no earlier head-to-head comparison of a weight edit with logit
bias for abstention; whether this edit improves on that baseline remains unresolved here.

\section{The experimental setup}
\label{sec:setup}

We evaluated Qwen2.5-3B-Instruct, a 36-layer model, in fp32 with greedy decoding. The prompt asks it to write
a short memory note and then an \texttt{Answer:} line containing a city or ``unknown''. The instruction
explicitly offers ``unknown'' when the memory does not supply an answer.

\paragraph{What the note tells us.}
The note is a prompted summary of the memory, not a verified reasoning trace. When we describe a response as
``note, then refuse'', we mean that the gold city appears in that summary before \texttt{Answer: unknown}.
This shows that the city was written down. It does not establish that the model completed an internal
derivation, and we do not treat the note as a faithful explanation of the computation.

\paragraph{Tasks and examples.}
The synthetic memories support three tasks: \textbf{direct extraction (DE)}, such as ``Where does A live?'';
\textbf{relational composition (RC)}, such as ``In which city does A's supervisor work?''; and
\textbf{temporal updating (TU)}, such as ``Where does A live now?'' after a move.
Each underlying \emph{base case} produces 13 related rows: 6 answerable rows (two variants at three hardness
levels), 3 unanswerable rows with a related distractor (UNA-rel), 1 unanswerable row with a swapped name
(UNA-swap), and 3 counterfactual rows in which the evidence specifies a different city (CF).

We call the unedited model \textbf{B0}. Its responses define the populations used throughout the evaluation:
\textbf{FA} contains answerable rows on which B0 abstains falsely; \textbf{C} contains answerable rows it gets
correct; \textbf{UNA} contains all unanswerable rows; and \textbf{CF} contains the changed-evidence rows.
Each intervention is evaluated on the same fixed populations.

\paragraph{Splits.}
Related variants stay together when cases are assigned to exploration, selection, confirmation and transfer
splits (\cref{tab:splits}). The partition also prevents a person--evidence-city key, including a
counterfactual variant's key, from crossing splits. No canonical base key or exact memory text appears in two
splits. Names, cities and templates can still be shared; disjoint cases do not imply entirely different
vocabularies. The agent saw only the exploration split, E.

\begin{table}[ht]
\centering\small
\begin{tabularx}{\linewidth}{@{}llYY@{}}
\toprule
Split & Bases per task & Access & Purpose \\
\midrule
E & 48 & Agent & Exploration only \\
S & 32 & Analysts & Dose selection and candidate selection \\
H & 48 & Sealed; read once & Confirmation \\
T-city & 24 & Sealed; read once & Transfer to new cities \\
\bottomrule
\end{tabularx}
\caption{E and S use generator seed 51 with disjoint cases, H uses seed 52, and T-city uses seed 53.
Rows from earlier experiments are excluded by prompt hash. A second planned transfer split with a reworded
instruction, T-format, was dropped before the first design freeze (F0).}
\label{tab:splits}
\end{table}

\paragraph{Scoring.}
A frozen rule-based grader scores the answer after the last \texttt{Answer:} marker, separately from the note
(\cref{app:methods}). It distinguishes a correct answer, abstention, wrong answer, mixed answer and abstention,
conflicting answers, other content, and degenerate or missing outputs.

The primary utility, $U$, rewards a correct answer with $+1$ and assigns abstention $0$. A mixed answer receives
$0$ on an answerable row and $-1$ on an unanswerable row; other errors receive $-1$. Within each base,
answerable and unanswerable rows each contribute half the weight. Counterfactual rows are scored separately
and do not enter $U$. We also report:
\begin{itemize}
  \item \textbf{Recovery (Rel):} the share of baseline false abstentions that become correct answers.
  \item \textbf{Unsupported answering (Uns):} the share of unanswerable rows given an asserted or suggested answer.
  \item \textbf{Damage (Dam):} the share of baseline-correct rows that are no longer correct.
  \item \textbf{Counterfactual following (CF-follow):} the share of changed-evidence rows answered with the new city.
  \item \textbf{Degenerate outputs (Deg):} degenerate responses and responses missing an answer marker.
\end{itemize}
Rates are base-weighted. We also give raw row counts where useful; these need not yield the same percentages.

\paragraph{Comparators.}
Five interventions serve as repair baselines, with settings frozen after selection on S:
\begin{itemize}
  \item \textbf{B1, logit bias:} penalise ``unknown'' tokens after the answer marker.
  \item \textbf{B2, added instruction:} append ``Do not reply unknown.'' to the prompt.
  \item \textbf{B3, option removal:} remove the ``unknown'' option from the instruction.
  \item \textbf{B4, activation steering:} apply the earlier refusal-direction push at layer 29.
  \item \textbf{B5, directional orthogonalisation:} remove a direction fitted on E from the relevant weights,
        following an adaptation of \citet{arditi2024refusal}.
\end{itemize}
A sixth comparator, \textbf{B6}, consists of five size-matched random edits. It is used in the pathway tests,
not in the main repair comparison (\cref{sec:explain}).

\begin{figure}[ht]
\centering
\begin{tabularx}{\linewidth}{@{}YYYY@{}}
\toprule
\textbf{Explore} & \textbf{Select} & \textbf{Freeze} & \textbf{Evaluate} \\
\midrule
Agent proposes edits on E; the harness replays them. &
S selects candidates and doses separately for repair and pathway tests. &
F2 fixes candidates, settings and analysis before confirmation. &
H tests the frozen claims; T-city checks transfer descriptively. \\
\bottomrule
\end{tabularx}
\caption{The agent sees only E. S selects the candidates and doses for repair and pathway tests. The
confirmatory plan is frozen by commit before any model outputs on H or T are generated; each is then read once.}
\label{fig:pipeline}
\end{figure}

\section{What the agent proposed}
\label{sec:proposals}

The editor was a fresh Claude Opus 5.5 session run through the API. It worked in a sandbox with read-only access
to the exploration split, E, and no access to this project's records or to S, H or T. Its tools could run the
model, record activations and apply candidate edits. Analysis was limited to closed-form operations such as
means, projections, SVD and single linear solves; gradient-based optimisation was not allowed.

An edit could affect at most four components: attention output projections (\texttt{o\_proj}), MLP output
projections (\texttt{down\_proj}) or layer norms. Allowed edits were rank-$\le$4 updates or sparse column edits.
The embedding, unembedding and rows for the ``unknown'' token were off limits. Before testing an edit, the
agent had to identify the component, state a hypothesis, predict outcomes for each row type, and describe
expected activity in the unedited model.

One run produced six saved edits and three submissions (\cref{tab:submissions}). It took 73 turns,
16 forward-pass requests and 63 Python runs, over about 16.4 hours, at an API cost of \$3.58.

\begin{table}[ht]
\centering\small
\begin{tabularx}{\linewidth}{@{}l>{\raggedright\arraybackslash}p{3.2cm}Y@{}}
\toprule
\# & Edit & The agent's hypothesis (abridged; not an established explanation) \\
\midrule
1 & Zero the \texttt{o\_proj} columns of heads L32h2 and L33h10 (repair) &
The two heads write a nearly constant ``unknown'' vector at the answer position, independent of the evidence.
Other late heads write the evidence city. False abstention occurs when the city signal is too weak to overcome
this default. Removing the default should recover false abstentions while preserving abstention on
unanswerable questions, where the ``unknown'' margin is larger. \\
2 & Rank-1 update to the layer-35 \texttt{o\_proj} (repair) &
Layer-35 attention contains an ``answer is available'' signal that is nearly as strong on false abstentions as
on correct answers, and absent on unanswerable questions. The edit reads that signal and writes toward city
tokens in the final layer, where no later layer can undo it. \\
3 & Double the same two heads (diagnostic) &
If the heads supply an abstention default, strengthening them should turn correct answers into ``unknown'',
even when the note names the city. \\
\bottomrule
\end{tabularx}
\caption{The agent's three submissions. Edits 1 and 3 change about 0.52M weights; edit 2 changes about 4.19M.
\todo{Check these abridged hypotheses against the transcript and add the key quotations.}}
\label{tab:submissions}
\end{table}

For edit 1, the agent predicted recovery of about 85\% of false abstentions, retention of at least 97\% of
correct answers, and an increase in unsupported answers on UNA-rel from about 1\% to at most 4\%.
For edit 2, it predicted 75--80\% recovery with essentially unchanged unsupported answering. For edit 3, it
predicted that 30--40\% of correct answers would become ``unknown''. These predictions were recorded for
evaluation; they did not determine candidate selection.

\section{The independent evaluation}
\label{sec:evaluation}

The evaluation separates two questions: whether an edit \emph{repairs} false abstention (branch R), and whether
the proposed \emph{pathway and explanation} survive intervention tests (branch M). Each branch can select a
different edit and dose. Success in one does not imply success in the other.

\paragraph{Replay and screening on E.}
After the agent finished, the evaluation harness replayed every submission on all of E. The repair screen
required a utility gain over B0 with a one-sided 97.5\% lower bound above zero and at least 20 false-abstention
bases in DE or RC. The pathway screen required an effect of at least 0.05 in the declared direction, with a
one-sided 97.5\% bound beyond zero and at least 20 bases in the relevant endpoint population. If no submission
met either screen, S, H and T would have remained unread. E outcomes remained exploratory because the agent
had used that split to develop the edits.

\paragraph{Selection on S.}
The entry rules were then evaluated on S. For repair, candidate and comparator doses were chosen to maximise
utility, and the qualifying candidate with the highest utility was selected; candidate ties favoured fewer
edited weights. For the pathway branch, dose selection maximised the declared signed effect, subject to the
pre-registered limit on degenerate outputs. Candidate selection used the standardised effect, separately from
repair. Dose ties favoured the setting closest to no change.

The repair candidate was submission 2, a rank-1 edit to the layer-35 attention output projection, at dose 2.
Its utility on S was 0.429, compared with 0.402 for the next submission and 0.303 for B0. Dose 2 was the top of
the frozen grid; larger doses were not searched. The pathway candidate was submission 1, the two-head edit,
at dose 1.5, with the layer-33 head designated as its primary component. Selected repair-comparator doses
were $-8$ for logit bias (B1), 1 for steering (B4), and $\beta = 1$ for orthogonalisation (B5).

\paragraph{Freezing the confirmation.}
The selected candidates, doses, comparator settings and analysis code were frozen before any H or T model
outputs were generated or read. This freeze, F2, is recorded at commit \texttt{02bb2f6}, with design hash
\texttt{d02d9f9c\ldots}. The confirmation run checked that commit before accessing H or T.

\paragraph{Decision criteria.}
H and T were each evaluated once. The repair criterion on H required a positive lower bound for the utility
gain, damage no greater than 0.05, degenerate outputs no greater than 0.02, and counterfactual following of at
least 0.90. The damage and counterfactual thresholds applied to point estimates; uncertainty was reported
alongside them.

A better trade-off than a repair comparator could be established in either of two ways:
\textbf{O1}, higher recovery with an unsupported-answer increase smaller than 0.05; or
\textbf{O2}, fewer unsupported answers with a recovery loss smaller than 0.05. Superiority and non-inferiority
had to pass together for the relevant orientation. Holm correction covered all ten comparisons---two
orientations for each of B1--B5. Pathway tests formed a separate family (\cref{sec:explain}).

\paragraph{Uncertainty.}
We used a paired bootstrap clustered by base and stratified by task, with $B = 10{,}000$ draws.
Intervals are 95\% unless stated otherwise. The intervals describe uncertainty in individual estimates;
confirmatory decisions use the corrected tests. \Cref{app:methods} gives the inference rules.

\section{Results and costs}
\label{sec:results}

% Main-text confirmation results come from H; E results remain in Appendix A.
The selected layer-35 edit met all four pre-registered repair criteria on H. Utility improved by
$+0.117$ (95\% CI 0.096 to 0.138). Damage to already-correct answers was 0.030, below the 0.05 limit on the
point estimate, although its upper bound was 0.055. Degenerate outputs were 0.008, below the 0.02 limit,
and counterfactual following was 0.917, above the 0.90 threshold.

\paragraph{Comparison with simpler interventions.}
The edit improved the release-versus-unsupported-answer trade-off over the added instruction (B2),
activation steering (B4) and directional orthogonalisation (B5). Their Holm-adjusted $p$-values were
0.005, 0.033 and 0.007, respectively (\cref{tab:h}).

The logit-bias comparison (B1) was inconclusive. The edit had higher recovery (0.816 versus 0.793) and fewer
unsupported answers (0.042 versus 0.095), but the comparison did not survive Holm correction
(adjusted $p = 0.092$). Equivalence was not established either: the two one-sided tests (TOST) gave
$p = 0.21$ for recovery and 0.70 for unsupported answers. The edit also caused damage that logit bias did not:
11 of 593 previously correct answers became incorrect, a base-weighted damage rate of 0.030.
Both logit bias and the layer-29 steering intervention damaged none of those answers.

Removing the ``unknown'' option (B3) recovered almost every false abstention, but produced unsupported answers
on 68\% of unanswerable questions. Superiority over this comparator was not established under the frozen test.

\begin{table}[ht]
\centering\small
\begin{tabularx}{\linewidth}{@{}Yrrrrrl@{}}
\toprule
Arm (on H) & Rel & Uns & Dam & CF-follow & $U$ & Edit better? \\
\midrule
B0 unedited & 0 & 0.035 & 0 & 0.632 & 0.317 & --- \\
\textbf{Selected edit} & \textbf{0.816} & \textbf{0.042} & \textbf{0.030} & \textbf{0.917} & \textbf{0.435} & --- \\
B1 logit bias ($-8$) & 0.793 & 0.095 & 0 & 0.919 & 0.406 & inconclusive \\
B2 ``do not reply unknown'' & 0.721 & 0.083 & 0.046 & 0.882 & 0.389 & yes (O2, 0.005) \\
B3 no ``unknown'' option & 0.977 & 0.682 & 0.021 & 0.975 & 0.130 & not established \\
B4 layer-29 push ($\alpha = 1$) & 0.786 & 0.078 & 0 & 0.921 & 0.418 & yes (O2, 0.033) \\
B5 orthogonalisation ($\beta = 1$) & 0.678 & 0.031 & 0.099 & 0.838 & 0.401 & yes (O1, 0.007) \\
\bottomrule
\end{tabularx}
\caption{Sealed confirmation on H: 48 bases per task, with 261 false-abstention (FA), 593 already-correct (C)
and 576 unanswerable (UNA) rows. Rel is recovery, Uns unsupported answering, Dam damage, and CF-follow
following changed evidence. Rates are base-weighted means, so they need not equal pooled row-count fractions;
$U$ uses $\lambda = 1$. The final column asks whether the selected edit has a better trade-off than that
comparator, giving the test and its Holm-adjusted $p$ across ten tests. Values were independently recounted
from per-row outputs. B6, the random-edit control, is used only in the pathway tests (\cref{sec:explain}).}
\label{tab:h}
\end{table}

% Planned figure: recovery versus unsupported answering for the selected edit and B1--B5, with CIs.

\paragraph{Transfer to new cities.}
T-city contains 24 bases per task with unseen cities. It was read once and provides a descriptive transfer
check, not a second confirmation. Results were consistent with H: utility was 0.422 versus 0.303 for the
unedited model, recovery was 0.817, unsupported answering was 0.062 versus 0.059, damage was 0.031,
and counterfactual following was 0.912. Comparator utilities retained their order:
B4 0.415, B1 0.412, B5 0.409, B2 0.397 and B3 0.107.
Transfer to a reworded instruction was not tested (\cref{sec:setup}).

\paragraph{Costs.}
The agent run cost \$3.58 in API usage. GPU evaluation cost \$4.38 for selection, including \$0.09 for a failed
first attempt, and \$2.57 for confirmation over 4.36 hours on one A40: \$6.95 in total GPU charges.
The exploration-split replay took about 59 hours on a laptop CPU, with \$0 in cloud charges.

\section{What the edits explain}
\label{sec:explain}

An edit can improve answers without explaining why the original model failed. We therefore distinguish
\emph{repair}, a beneficial change in behaviour; \emph{pathway evidence}, an intervention linking the edit's
effect to a named component; and a \emph{mechanistic account}, a description that also predicts the unedited
model's behaviour. The repair result above concerns the layer-35 edit. The pathway evaluation selected a
different candidate: the two-head edit, at dose 1.5, with the layer-33 head fixed in advance as the primary
component.

\paragraph{A pathway through layer 33.}
Both pathway tests passed after Holm correction across the three branch-M tests ($p < 0.001$).
Reverting the layer-33 component in the edited model reduced recovery on false abstentions by 0.293
(95\% CI 0.225 to 0.368). The edit also recovered 0.717 more than the mean of five size-matched random edits,
B6. These tests address the effect of changing that component and the comparison with random perturbations.

Two activation interventions asked whether the effect could be transferred or blocked. Adding the edit's
change to the layer-33 output at the answer position in the unedited model recovered 0.808 of false
abstentions. Clamping that output back to its original value in the edited model removed only 0.193.
The layer-33 intervention was therefore sufficient, but not necessary, to release answers under these
conditions. The remaining recovery is consistent with another route, possibly through the layer-32 head;
that explanation was not directly tested. Without a random-head comparison, these results do not establish
that the two selected heads have a unique role.

\paragraph{The explanation was not established.}
The agent's account of what the component computes was not supported by the frozen tests. Its proposed
variable added no predictive value for false abstention, although the baseline predictor already had an
AUROC of 0.998 and left little room for improvement. Intervening on the variable also failed to establish a
different effect on answerable versus unanswerable questions relative to a matched logit bias: the
interaction was 0.0005 (95\% CI $-0.025$ to 0.027). Two activity predictions held, but the combined account
test did not pass.

This negative has an important qualification. The agent described the heads as an evidence-independent
``unknown'' prior. A null interaction is close to what that description would predict, so the frozen predictor
and interaction tests fit this account poorly. The result means that the account was not shown by the
pre-registered tests; it is not a refutation of that account.

At its pathway-test dose, the two-head edit was also not selective: on S it produced unsupported answers on
96\% of unanswerable rows. The pathway result identifies where an answer-release effect can be produced.
It does not explain the selected layer-35 repair.

\paragraph{Exploratory observations.}
On E, doubling the two heads changed 180 of 572 correct answers (31\%) into ``unknown'', within the agent's
prediction of 30--40\%. Removing the heads recovered most false abstentions on that split. These observations
motivated the bidirectional interpretation, but E was the split the agent studied and cannot independently
confirm it. The exploratory results are collected in \cref{app:e}.

\section{Follow-up: are the two edits specific?}
\label{sec:followup}

% Design text: proposal rev 1b, manifest b2b14bac, committed a177536 before held output.
% Leave results pending until the Advisor recount is complete.

The confirmation leaves open whether the selected components and directions matter, or whether other edits
of the same kind would work as well. A follow-up was pre-registered after confirmation closed. It uses a new
controlled split, N (seed 54), disjoint from the agent's exploration and the selection split. The follow-up
addresses three questions without changing the confirmation verdicts.

\paragraph{Does the read direction matter?}
The agent described the layer-35 edit as reading an ``answer available'' signal and writing along
city~$-$~``unknown''. Twenty control edits keep the write vector but use a random read vector. Each is scaled
on development rows to match the edit's root-mean-square gate and sign-aligned to push in the same direction.
The frozen criterion, C4, requires the edit's utility to exceed all twenty controls and its unsupported-answer
rate to fall below their median. Failure would leave read-side specificity unestablished; it would not by
itself show equivalence to an output bias.

\paragraph{Are the two heads unusual among late-layer heads?}
C1 requires the target pair's ablation to recover more false abstentions than each of twenty random pairs
from the same layers, with a margin of at least 0.30 over their median. C2 requires doubling the target
pair to create more abstentions on previously correct rows than doubling each random pair.

\paragraph{Does the layer-35 signal follow the agent's predictions?}
C5 tests four frozen predictions about changes in the proposed signal when evidence is removed, the city
changes, distractors are added, or unrelated evidence replaces absent evidence.

C1, C2 and C4 use an empirical rank of 1 among 21 evaluated edits; C5 uses an intersection--union test at
0.05. No family-wise error-control claim is made. Under a null model in which target scores are exchangeable
with the random-control scores, with ties counted against the target, and with a valid C5 test, a union bound
puts the chance of at least one passing claim at most about 0.19. This error bound applies under the stated
exchangeability null and a valid C5 test; it is not a guarantee for every possible mechanism of false abstention.

Descriptive arms examine single heads, other late-layer pairs, random write vectors, direct logit attribution,
and the prompt with its ``unknown'' option removed. A natural-question check uses 138 LongMemEval development
questions with evidence in view.

\paragraph{Results.} \pending{§53 results: C1, C2, C4, C5 verdicts; LongMemEval releases, damage and answers on
abstention questions; spend.}

\section{Limits and reproducibility}
\label{sec:limits}

\paragraph{Scope.}
This study concerns one model, Qwen2.5-3B-Instruct, on synthetic city questions. The new-city transfer set
checks a change in answer vocabulary, not a new task or prompt format. T-format was dropped before F0, so
transfer to reworded instructions is untested. Natural-question evaluation on LongMemEval development rows
is part of the follow-up (\cref{sec:followup}); no such result is claimed here. Misleading-context conditions
such as those in GRAB-RAG \citep{setiawan2026grabrag} would provide a further test of evidence use.

\paragraph{What the controls leave open.}
The only random-edit comparator in confirmation was B6, five size-matched perturbations used in the pathway
test. Other head pairs and rank-1 edits with random read directions are evaluated in the follow-up.
The prompt offers ``unknown'', and the present results do not distinguish a change in how the model handles
that option from a change in how it uses evidence.

The frozen comparator set also omits component-level methods such as Steer2Edit \citep{steer2edit2026},
weight-folded false-refusal ablation \citep{wang2025falserefusal}, gated refusal edits
\citep{hinkley2026paving}, and SAE-feature steering for abstention \citep{hasanov2026drno}.
Our comparisons therefore do not establish superiority over editing or steering methods in general.

\paragraph{Precision and costs to performance.}
The logit-bias comparison is unresolved. The recovery difference has a 95\% CI of $-0.043$ to $+0.088$;
superiority and equivalence were not established at the frozen margins. Counterfactual following, 0.917,
cleared its 0.90 threshold by a small margin. The selected edit damaged previously correct answers at a
base-weighted rate of 3\%, while B1 and B4 damaged none. Meeting the repair criterion does not remove that cost.

The selected dose was the highest in the frozen grid, and utility was still increasing on S. Larger doses
were not searched, so the selected edit is not claimed to be the best available version of its kind.
The pathway account test also faced a predictor ceiling, with baseline AUROC 0.998; its negative must be
read alongside the interaction test and the account--test mismatch described in \cref{sec:explain}.

\paragraph{One agent run.}
Other sessions or agents may propose different edits. We did not test whether access to weights was necessary
for this agent's success. In a related auditing setting, a calculation tool mattered more than weight access
\citep{liu2026auditors}. The exploration split E was visible to the agent, so its results remain exploratory
and are collected in \cref{app:e}.

\paragraph{Reproducibility.}
The design was frozen in three stages, F0, F1 and F2. The final freeze preceded all model outputs on H and
T-city. Code, edit recipes, the agent transcript and per-row outputs are listed in \cref{app:artefacts}.


\section*{Disclosure}
The weight edits were proposed by Claude Opus 5.5 (Anthropic), run through the API. AI assistants (Claude and
GPT models) also helped design the study, write analysis code, review it, and draft this report. Every number in
the main text was recounted from per-row outputs. \todo{final wording.}


\bibliographystyle{plainnat}
\bibliography{refs}

\appendix
\section{Exploratory results on the agent's split (E)}
\label{app:e}

These numbers come from E, the split the agent studied while choosing its edits. They show which submissions
entered the evaluation; they are not evidence that the edits generalise. E has 144 bases (48 per task) and
1,872 rows; B0 abstains falsely on 284 answerable rows (94 bases: DE 44, RC 35, TU 15).

\begin{table}[ht]
\centering\small
\begin{tabularx}{\linewidth}{@{}lYYY@{}}
\toprule
Entry on E & Edit 1 (zero heads) & Edit 2 (rank-1 L35) & Edit 3 (double heads) \\
\midrule
R: $U - U_{B0}$ & $+0.122$ ($+0.102$, $+0.141$) & $+0.122$ ($+0.104$, $+0.139$) & $-0.107$ ($-0.126$, $-0.088$) \\
M: declared effect & $+0.817$ ($+0.753$, $+0.877$) & $+0.706$ ($+0.633$, $+0.777$) & $+0.360$ ($+0.308$, $+0.410$) \\
Enters & R and M & R and M & M only \\
\bottomrule
\end{tabularx}
\caption{Entry on E. Intervals are one-sided 97.5\% bootstrap bounds. The M effect is FA$\rightarrow$gold for
edits 1 and 2 and C$\rightarrow$abstention for edit 3, from the controlled prefix.}
\end{table}

\begin{table}[ht]
\centering\small
\begin{tabular}{@{}lcccc@{}}
\toprule
Rate on E (base-weighted) & B0 & Edit 1 & Edit 2 & Edit 3 \\
\midrule
Rel (FA released to gold) & 0.000 & 0.862 & 0.739 & 0.009 \\
Uns (unsupported on UNA) & 0.040 & 0.076 & 0.050 & 0.030 \\
Dam (correct answers lost) & 0.000 & 0.028 & 0.015 & 0.381 \\
CF-follow & 0.653 & 0.940 & 0.887 & 0.451 \\
Deg & 0.005 & 0.006 & 0.003 & 0.013 \\
$U$ ($\lambda = 1$) & 0.305 & 0.427 & 0.427 & 0.198 \\
\bottomrule
\end{tabular}
\caption{Rates on E, total effect.}
\end{table}

\paragraph{Row counts.}
Edit 1 released 252 of 284 false abstentions and kept 561 of 572 correct answers; unsupported answers rose from
6 to 15 of 432 UNA-rel rows and from 17 to 29 of 144 UNA-swap rows. Edit 2 released 222 of 284 and kept 563 of
572; unsupported answers rose from 6 to 11 and from 17 to 18. Edit 3 turned 180 of 572 correct answers into
``unknown'' (the agent predicted 30--40\%).

\section{Claims and evidence}
\label{app:claims}

\begin{table}[ht]
\centering\small
\begin{tabularx}{\linewidth}{@{}YY>{\raggedright\arraybackslash}p{0.23\linewidth}@{}}
\toprule
Claim & Evidence & Status \\
\midrule
False abstentions in the 3B are mostly ``note, then refuse'' & 152 of 165 held false abstentions name the gold in the note (earlier experiment) & established \\
Earlier ungated pushes recover answers but also increase unsupported answering & earlier held results (Uns 0.188 to 0.451); a separate template-limited gating result is discussed in the text & established \\
At least one agent-proposed edit enters the evaluation & entry on E (\cref{app:e}) & exploratory \\
The selected edit repairs false abstention on sealed data & H, repair criterion (\cref{sec:evaluation}) & shown \\
It has a better trade-off than B2, B4 and B5 & H, Holm-corrected O1/O2 tests & shown \\
It has a better trade-off than B1 or B3 & H, Holm-corrected O1/O2 tests; TOST for B1 & not established \\
It carries over to new cities & T-city & consistent (descriptive) \\
The two-head edit's release runs through the layer-33 head & H, branch M tests M1 and M2 & sufficient under the tested intervention; not necessary \\
The agent's account of what that component computes & H, branch M test M3 & not supported \\
The layer-35 edit's read side matters (C4) & N, follow-up (\cref{sec:followup}) & \pending{§53} \\
The two heads are special among late-layer pairs (C1, C2) & N, follow-up & \pending{§53} \\
The agent's description of the layer-35 signal (C5) & N, follow-up & \pending{§53} \\
The edits carry over to natural questions & LongMemEval dev, follow-up & \pending{§53} (descriptive) \\
\bottomrule
\end{tabularx}
\caption{Every claim in the main text, with the evidence that supports it.}
\end{table}

\section{Artefacts}
\label{app:artefacts}

\begin{itemize}
  \item Design freezes: F0 \texttt{d29263a}; F1 \texttt{1a37a72} (built at \texttt{c9ca458}); F2 \texttt{02bb2f6} (design hash \texttt{d02d9f9c\ldots}).
        Pod tooling fixes: \texttt{b6e3d8e} (before any pod); \texttt{c880239} and \texttt{bf22ac3} (before the
        confirmation pod). Confirmation outputs: \texttt{321a4f1} (report sha256 \texttt{51bbbd2e\ldots}).
  \item Case file: 5,928 rows, 456 bases, sha256 \texttt{5eeb7c40\ldots}.
  \item Edit recipes: \texttt{c64ea090781450cf}, \texttt{8364bf135ac59a83}, \texttt{87224655b7a5361e}
        (head columns and rank-1 vectors). Edited weights are released only if the model licence allows.
  \item Follow-up (\cref{sec:followup}): proposal rev 1b \texttt{ea88023}; build \texttt{9ff3d61} and sign rule
        \texttt{1dce499}; manifest sha256 \texttt{b2b14bac\ldots}, committed at \texttt{a177536} before any held
        output; N split 1,950 rows, sha256 \texttt{651c930c\ldots}. \pending{§53 output commit and report sha.}
  \item The agent's full transcript, with paths and keys removed.
  \item Per-row outputs for every arm and split.
  \item \todo{public repository URL; Zenodo DOI if used.}
\end{itemize}

\section{Related work in detail}
\label{app:related}

\begingroup\small
\begin{longtable}{@{}>{\raggedright\arraybackslash}p{0.25\linewidth}>{\raggedright\arraybackslash}p{0.34\linewidth}>{\raggedright\arraybackslash}p{0.35\linewidth}@{}}
\caption{Closest work. Descriptions and bibliographic details marked for verification need a primary-source check before release.}\label{tab:related}\\
\toprule
Work & What it did & How this report differs \\
\midrule
\endfirsthead
\toprule
Work & What it did & How this report differs \\
\midrule
\endhead
\bottomrule
\endfoot
Weight-space programming \citep{takehi2026wsp} & an LLM agent edits a small fraction of Qwen3.5-0.8B's weights to stop reasoning loops, with held-out evaluation & a behaviour (looping) with no correct refusals to protect; here the score penalises new guessing on unanswerable questions, and confirmation is sealed and pre-registered \\
Steer2Edit \citep{steer2edit2026} & steering vectors turned into rank-1 edits of individual heads and neurons & not one of our frozen comparators; future work \\
Safety heads \citep{zhou2025safetyheads} & ablating one attention head in Llama-2-7b-chat sharply raises compliance with harmful requests & the refusal there is wanted; ours is an unwanted ``unknown'' with the evidence in view \\
Retrieval heads \citep{wu2025retrievalheads} & under 5\% of heads do in-context retrieval; pruning them causes hallucination, pruning random heads does not & our interventions target false abstention with sufficient evidence; the role of the selected heads remains unresolved. Random-head controls are part of our follow-up \\
Refusal and unanswerability directions \citep{arditi2024refusal,lavi2025unanswerability,lee2025cast} & a single direction mediates refusal; an unanswerability direction steers abstention; conditional steering & our evaluated ungated pushes also increased unsupported answering; an earlier gate was template-limited \\
False-refusal vector ablation \citep{wang2025falserefusal} & a false-refusal direction folded into the weights; scored on released false refusals against kept true refusals & safety over-refusal, not evidence in view; not agent-proposed; not one of our frozen comparators \\
Residual Paving \citep{hinkley2026paving} & gated residual-stream edits of safety refusals with a three-way selectivity scheme, on a held-out split & the closest selectivity yardstick; safety refusal, no weight change, not agent-proposed, not pre-registered \\
Dr.\ No \citep{hasanov2026drno} & SAE-feature steering that balances answering against abstention in vision-language models & activation steering only; no weight edit, no pre-registration \\
GRAB-RAG \citep{setiawan2026grabrag} & benchmark of small frozen RAG models under supportive, missing and misleading context & no intervention; a ready source of natural counterfactual evidence for future tests \\
LLMs as auditors \citep{liu2026auditors} & an agent auditing small networks did as well without weight access when it had a calculation tool & auditing, not editing; a caution that an agent's account of weights may not drive its success, consistent with branch M \\
Pre-registered interpretability \citep{readablenotsteerable2026} & pre-registered, once-only held-out tests of representation steering & precedent for the method, outside language models \\
Abstention as prompt artefact; answer release \citep{ling2025artifact,liu2026broken} & abstention can follow the ``unknown'' option; answers are recoverable after refusal and release is local & predicts that a late edit can release answers; this part is expected, not new \\
Commit--abstain circuit \citep{commitabstain2026} & a commit-versus-abstain circuit across ten models & different models and setting; no factual grading of released answers \\
Abstention behaviour \citep{joren2025sufficient,kirichenko2025abstentionbench,kumaran2026confidence} & models misuse sufficient context; over- and under-abstention; confidence thresholds drive abstention & behavioural or closed-book accounts; no component-level localisation of in-view false abstention \\
Interpretability agents \citep{shaham2024maia,barez2026agenda} & agents that describe components or propose interventions & here one agent-proposed intervention is tested against pre-registered bars \\
Knowledge editing \citep{meng2022rome,meng2023memit} & locate and edit stored facts & a decision behaviour, not a stored fact \\
\bottomrule
\end{longtable}
\endgroup

\section{Methods in detail}
\label{app:methods}

\paragraph{Answer span and grading.}
The grader uses the text after the last \texttt{Answer:} marker and scores the preceding note separately.
A missing marker is its own class and contributes to the reported degenerate-output rate.
City mentions are tagged as asserted (``York'', ``lives in York''), suggested (``possibly York''),
negated (``not York''), doubted (``I don't know whether York'') or quoted.
Only asserted and suggested mentions count as answers. With no asserted or suggested answer, an abstention
phrase yields the abstention class. One city is graded as gold, wrong or other content; one city together
with an abstention phrase is mixed; two or more distinct asserted or suggested cities are conflicting.
The grader and its fixtures were frozen before the agent saw E.

\paragraph{Utility.}
\begin{center}\small
\begin{tabularx}{\linewidth}{@{}lcccY@{}}
\toprule
Row & Gold & Abstention & Mixed & Wrong, conflicting, other, degenerate or missing \\
\midrule
Answerable, counterfactual & $+1$ & $0$ & $0$ & $-\lambda$ \\
Unanswerable & --- & $0$ & $-\lambda$ & $-\lambda$ \\
\bottomrule
\end{tabularx}
\end{center}
$\lambda = 1$ is primary; $\lambda = 2$ is a secondary setting. Each base contributes equally to aggregate
utility. Within a base, answerable and unanswerable rows each supply half the weight. Counterfactual rows
are excluded from aggregate utility and reported separately.

\paragraph{Two intervention settings.}
For the \emph{total effect}, the model generates both its note and its answer under the intervention. This
is the primary setting for repair. For the \emph{controlled-prefix effect}, the intervention recomputes the
unedited model's stored token IDs through the answer marker, then generates only the answer. This is the
primary setting for pathway tests. Its shared text prefix does not imply a shared internal state.

\paragraph{Inference.}
The paired percentile bootstrap is clustered by base and stratified by task, with $B = 10{,}000$ attempted
draws and a fixed seed. Tests use approximate percentile inversion. For a one-sided null
$\theta \leq \theta_0$, the tail score is
\[
  p = \frac{1 + \#\{b : \theta_b^* \leq \theta_0\}}{B_{\mathrm{valid}} + 1}.
\]
Undefined replicates, such as those with an empty denominator or undefined AUROC, are excluded and counted.
If more than 1\% of attempted draws fail, the statistic receives no decision. If all valid draws equal the
point estimate, the frozen rule assigns $1/(B_{\mathrm{valid}}+1)$ when that estimate is on the alternative
side of the threshold, and 1 otherwise. This degenerate-distribution rule differs from dropping undefined
draws. These are approximate bootstrap tests, not exact randomization tests.

Equivalence uses two one-sided tests (TOST). For an intersection--union claim, the composite $p$ is the
largest elementary $p$. The repair family contains O1 and O2 against B1--B5, ten hypotheses; the pathway
family contains M1--M3. Each family is Holm-corrected at 0.05, with no study-wide guarantee claimed.
Marginal 95\% intervals describe uncertainty; corrected tests determine the confirmatory decisions.
The 0.05 margins, 0.90 counterfactual-following threshold and 20-base minimum are operational conventions.

\paragraph{Stopping rules.}
H and T were evaluated once, with no rescue variant or second agent run. If no candidate met the E screening
rule or the S entry rule, the study would have stopped and documented that outcome without opening H or T.


\end{document}
